\documentclass[11pt]{article}
\usepackage[margin=1in]{geometry}
\usepackage{amsmath,amssymb,amsthm,mathtools}
\usepackage[hidelinks]{hyperref}

\newtheorem{theorem}{Theorem}[section]
\newtheorem{lemma}[theorem]{Lemma}
\newtheorem{corollary}[theorem]{Corollary}
\newtheorem{proposition}[theorem]{Proposition}
\theoremstyle{definition}
\newtheorem{definition}[theorem]{Definition}
\theoremstyle{remark}
\newtheorem{remark}[theorem]{Remark}

\newcommand{\cl}{\operatorname{cl}}
\newcommand{\rk}{\operatorname{rk}}
\newcommand{\lht}{\operatorname{ht}}
\newcommand{\si}{\operatorname{si}}
\newcommand{\cL}{\mathcal L}
\newcommand{\cB}{\mathcal B}
\newcommand{\cC}{\mathcal C}
\newcommand{\cI}{\mathcal I}
\newcommand{\qbinom}[2]{\genfrac{[}{]}{0pt}{}{#1}{#2}_q}

\title{Boolean Antichains in Finite Lattices\\
Complete Classification, Product Structure, and Matroid Applications}
\author{Alex Chengyu Li}
\date{31 August 2026}
\hypersetup{
  pdftitle={Boolean Antichains in Finite Lattices: Complete Classification, Product Structure, and Matroid Applications},
  pdfauthor={Alex Chengyu Li}
}

\begin{document}
\maketitle

\begin{abstract}
Garber, Goltermann, Horiatakis, K\"onig, and Gottesman asked for
constructions and counts of Boolean antichains in geometric lattices and
proposed a spanning-tree correspondence for graphic matroids.  We solve these
problems in greater generality by classifying Boolean antichains of every size
in every finite lattice.  For an ordered family, intersecting all but one
member produces candidate Boolean atoms.  The family is Boolean precisely
when those atoms rise above the common meet and all $2^k$ reconstruction
height gaps vanish.  This gives an exact finite formula for the complete
enumerator.  Its binomial transform is multiplicative under lattice products,
equivalently giving an explicit positive convolution.

For a rank-$r$ matroid, maximum Boolean antichains in the flat lattice are
naturally in bijection with bases of the simplification.  For graphs this
gives spanning forests and specializes to the proposed spanning-tree
bijection in the connected simple case.  Retaining parallel classes recovers
the multivariate basis polynomial, while flat intervals give minor-local and
rank-tight versions.  The framework also yields explicit all-size formulas
for uniform matroids, complete classifications for finite distributive and
subspace lattices, a universal M\"obius formula for Boolean pairs, and
complete distributions for finite projective planes.
\end{abstract}

\section{Introduction}

Let $L$ be a finite lattice.  A Boolean antichain is an antichain $C$ for
which the map
\[
  S\longmapsto \bigwedge S
\]
is a lattice anti-isomorphism from the power set of $C$ to the sublattice
generated by $C$ and the greatest element of $L$.  This definition was
introduced for representation-theoretic purposes and was recently studied
enumeratively by Garber, Goltermann, Horiatakis, K\"onig, and Gottesman
\cite{GarberEtAl2026}.  Among other results, they proved that Boolean
antichains of size $n-1$ in the partition lattice are counted by the labelled
trees on $[n]$.  Since the partition lattice is the lattice of flats of the
graphic matroid of $K_n$, they expected the construction to extend to an
arbitrary graph and its spanning trees.  They also asked more generally for
constructions and counts in geometric lattices.

The main result is lattice-wide.  Given an ordered tuple
$\mathbf H=(H_1,\ldots,H_k)$, let $X_{\mathbf H}$ be its common meet and let
$A_i(\mathbf H)$ be the meet of all entries except $H_i$.  For
$S\subseteq[k]$, compare the join of the $A_i$ with $i\in S$ to the meet of
the complementary $H_j$.  The former always lies below the latter.  Their
height difference is the reconstruction gap $\gamma_{\mathbf H}(S)$.
Theorem~\ref{thm:global-gap} proves that $\mathbf H$ orders a Boolean
antichain precisely when every $A_i$ rises above $X_{\mathbf H}$ and every
reconstruction gap vanishes.  Consequently,
\[
 |\cC_k(L)|=\frac1{k!}
 \sum_{\mathbf H\in L^k}
 \prod_i\mathbf 1_{\{\lht_L A_i(\mathbf H)>\lht_L X_{\mathbf H}\}}
 \prod_{S\subseteq[k]}\mathbf 1_{\{\gamma_{\mathbf H}(S)=0\}}.
\]
Thus the complete enumerator is determined for every finite lattice, without
a distributivity, modularity, or grading hypothesis.

The second structural result describes products.  If
$c_k(L)=|\cC_k(L)|$ and
\[
 \eta_k(L)=\sum_{a=0}^k\binom ka a!c_a(L),
\]
then Theorem~\ref{thm:product} gives
\[
 \eta_k(L\times K)=\eta_k(L)\eta_k(K).
\]
Equivalently, the coefficients satisfy an explicit positive convolution.
This product law recovers the Stirling recurrence for Boolean lattices and
computes every layer of a matroid direct sum from its summands.

The motivating geometric-lattice problem is the principal application of
this framework.  Every geometric lattice is the lattice of flats of a simple
matroid \cite{Proudfoot2026}.  The subsets of a matroid basis are known to
span a Boolean sublattice \cite{Speyer2021}; the new point is the inverse
classification and its enumerative refinements.  If $M$ has rank $r$, the
atoms of any size-$r$ Boolean span are rank-one flats forming a basis of
$\si(M)$, and the construction is reversible.  Hence
\[
 \left\{\text{Boolean antichains of size }r\text{ in }\cL(M)\right\}
 \longleftrightarrow
 \left\{\text{bases of }\si(M)\right\}.
\]
For a graph, this is the spanning-forest bijection; the connected simple case
is exactly the spanning-tree extension proposed by Garber et al.

The flat lattice remembers each parallel class as a single atom.  Retaining
the class recovers the full multivariate basis polynomial:
\[
 \sum_{B\in\cB(M)}\prod_{e\in B}x_e
 =\sum_C\prod_{A\in\operatorname{At}(\langle C\rangle)}
       \left(\sum_{e\in A\setminus\cl(\varnothing)}x_e\right).
\]
The same statement holds in every flat interval through the minor
$(M|Y)/X$, yielding the rank-tight classification by contraction bases.  The
global height-gap theorem supplies all remaining layers, whose face ranks are
strict integral polymatroid ranks.

Several families admit more compact evaluations.  Boolean pairs have a
universal M\"obius formula; finite distributive lattices are described by
disjoint order filters; finite subspace lattices are described by quotient
direct-sum decompositions; and uniform matroids and finite projective planes
have explicit all-size formulas.  Section~2 introduces the motivating
flat-lattice notation.  Section~3 proves the global height-gap and product
theorems.  The remaining sections develop the universal specializations and
the matroid applications.

\section{Matroid motivation and flat-lattice notation}

Let $M$ be a finite matroid on ground set $E$, with rank function $\rk_M$ and
closure operator $\cl_M$.  Its lattice of flats is denoted by $\cL(M)$.  The
meet is intersection and the join is
\[
 F\vee G=\cl_M(F\cup G).
\]
The least flat is $L_0=\cl_M(\varnothing)$, the set of loops, and the greatest
flat is $E$.  A maximal chain adds one rank at each cover, so the height of
$\cL(M)$ is $r=\rk(M)$.

The simplification $\si(M)$ deletes the loops and replaces each nonloop
parallel class by a single element.  We identify its ground set with the
rank-one flats of $M$: a rank-one flat $A$ represents the parallel class
$A^\circ=A\setminus L_0$.  A basis of $\si(M)$ is therefore a set
$\{A_1,\ldots,A_r\}$ of rank-one flats whose join is $E$.

For a Boolean antichain $C=\{H_1,\ldots,H_k\}$, write
\[
 \langle C\rangle=\left\{\bigwedge_{i\in S}H_i:S\subseteq[k]\right\},
 \qquad \bigwedge\varnothing=E.
\]
Its bottom is $X_C=\bigcap_i H_i$.  The atoms of the Boolean lattice
$\langle C\rangle$ are
\begin{equation}\label{eq:boolean-atoms}
% formal-claim-begin: def-boolean-atoms
 A_i(C)=\bigcap_{j\ne i}H_j.
% formal-claim-end: def-boolean-atoms
\end{equation}

\begin{lemma}[Full-height rigidity]\label{lem:full-height}
% formal-claim-begin: lem-full-height
Let $L$ be a ranked lattice of rank $r$, and let $C$ be a Boolean antichain
of size $r$.  Then the bottom of $\langle C\rangle$ is the bottom of $L$,
and every maximal chain of $\langle C\rangle$ is a maximal chain of $L$.
In particular, the atoms and coatoms of $\langle C\rangle$ have ambient
ranks $1$ and $r-1$, respectively.
% formal-claim-end: lem-full-height
\end{lemma}

\begin{proof}
% relation-claim-begin: proof-lem-full-height
A maximal chain in $\langle C\rangle\cong B_r$ has $r$ cover steps.  Viewed
in $L$, each strict step raises rank by at least one.  Its top has rank $r$,
so its bottom has rank zero and every step raises rank by one.  A ranked
lattice has a unique element of rank zero, and the remaining assertions
follow.
% relation-claim-end: proof-lem-full-height
\end{proof}

We will also use the following familiar basis fact in a form that records
meets as well as joins.

\begin{lemma}[Boolean span of a simplified basis]\label{lem:basis-span}
% formal-claim-begin: lem-basis-span
Let $\mathcal A=\{A_1,\ldots,A_r\}$ be a basis of $\si(M)$.  For
$S\subseteq[r]$, put $F_S=\bigvee_{i\in S}A_i$, with $F_\varnothing=L_0$.
Then
\[
 F_S\vee F_T=F_{S\cup T},\qquad F_S\cap F_T=F_{S\cap T}.
\]
Consequently $S\mapsto F_S$ embeds $B_r$ as a sublattice of $\cL(M)$.
% formal-claim-end: lem-basis-span
\end{lemma}

\begin{proof}
% relation-claim-begin: proof-lem-basis-span
The join identity is immediate.  Choose one representative from each
$A_i^\circ$.  These representatives form a basis of $M$, so
$\rk(F_S)=|S|$.  Submodularity gives
\[
 \rk(F_S\cap F_T)
 \le |S|+|T|-|S\cup T|=|S\cap T|.
\]
The reverse inequality follows from $F_{S\cap T}\subseteq F_S\cap F_T$.
The two flats have the same rank and one contains the other, hence they are
equal.
% relation-claim-end: proof-lem-basis-span
\end{proof}

\section{Global height-gap classification}

Let $L$ now be an arbitrary finite lattice, and let $\lht_L(Z)$ be the maximum
length of a chain from its bottom to $Z$.  Strict order implies strict increase
of $\lht_L$.  The preceding definition of $A_i(C)$ makes sense for an arbitrary
ordered tuple $\mathbf H=(H_1,\ldots,H_k)\in L^k$, whether or not it is
Boolean.
Put
\[
 X_{\mathbf H}=\bigwedge_{i=1}^k H_i,
 \qquad
 A_i(\mathbf H)=\bigwedge_{j\ne i}H_j.
\]
For $S\subseteq[k]$, define two candidate faces
\begin{align*}
 U_{\mathbf H}(S)
   &=X_{\mathbf H}\vee\bigvee_{i\in S}A_i(\mathbf H),\\
 V_{\mathbf H}(S)
   &=\bigwedge_{j\notin S}H_j,
\end{align*}
where the empty relative join is $X_{\mathbf H}$ and the empty meet is
$\hat1$.
Every $A_i(\mathbf H)$ lies below every $H_j$ with $i\ne j$, and hence
\[
 U_{\mathbf H}(S)\le V_{\mathbf H}(S).
\]
The \emph{reconstruction height gap} of $S$ is the nonnegative integer
\begin{equation}\label{eq:reconstruction-gap}
% formal-claim-begin: def-reconstruction-gap
 \gamma_{\mathbf H}(S)
 =\lht_L V_{\mathbf H}(S)-\lht_L U_{\mathbf H}(S).
% formal-claim-end: def-reconstruction-gap
\end{equation}

\begin{theorem}[Global height-gap criterion]\label{thm:global-gap}
% formal-claim-begin: thm-global-gap-criterion
Let $L$ be a finite lattice and let
$\mathbf H=(H_1,\ldots,H_k)\in L^k$.  The following statements are
equivalent:
\begin{enumerate}
\item $\mathbf H$ is an ordering of a size-$k$ Boolean antichain;
\item the conditions
\begin{align}
 \lht_L A_i(\mathbf H)&>\lht_L X_{\mathbf H}
     &&\text{for every }i\in[k],\label{eq:atomic-rise}\\
 \gamma_{\mathbf H}(S)&=0
     &&\text{for every }S\subseteq[k].\label{eq:all-gaps-zero}
\end{align}
hold.
\end{enumerate}
% formal-claim-end: thm-global-gap-criterion
% relation-claim-begin: lead-global-enumerator
Consequently, writing $\mathbf 1_{\{P\}}$ for the indicator of a proposition
$P$, the exact count is
% relation-claim-end: lead-global-enumerator
\begin{equation}\label{eq:global-enumerator}
% formal-claim-begin: eq-global-enumerator
 \boxed{
 |\cC_k(L)|=\frac1{k!}
 \sum_{\mathbf H\in L^k}
 \prod_{i=1}^k
   \mathbf 1_{\{\lht_L A_i(\mathbf H)>\lht_L X_{\mathbf H}\}}
 \prod_{S\subseteq[k]}
   \mathbf 1_{\{\gamma_{\mathbf H}(S)=0\}}.}
% formal-claim-end: eq-global-enumerator
\end{equation}
% formal-claim-begin: thm-global-boundaries
This formula includes $|\cC_0(L)|=1$ and is zero for $k>\lht_L(\hat1)$.
% formal-claim-end: thm-global-boundaries
\end{theorem}

\begin{proof}
% relation-claim-begin: proof-global-criterion
Suppose first that $\mathbf H$ orders a Boolean antichain.  In its Boolean
span, $X_{\mathbf H}$ is the bottom and the $A_i(\mathbf H)$ are the atoms,
so Equation~\eqref{eq:atomic-rise} holds.  The face indexed by $S$ is both
the join of its atoms and the meet of the complementary coatoms.  Thus
$U_{\mathbf H}(S)=V_{\mathbf H}(S)$ and every reconstruction gap vanishes.

Conversely, assume Equations~\eqref{eq:atomic-rise} and
\eqref{eq:all-gaps-zero}.  Since $U_{\mathbf H}(S)\le V_{\mathbf H}(S)$,
equality of their heights implies equality of the two elements.  Write their
common value as $F_S$.  Then
\begin{align*}
 F_S\wedge F_T
 &=V_{\mathbf H}(S)\wedge V_{\mathbf H}(T)
   =V_{\mathbf H}(S\cap T)=F_{S\cap T},\\
 F_S\vee F_T
 &=U_{\mathbf H}(S)\vee U_{\mathbf H}(T)
   =U_{\mathbf H}(S\cup T)=F_{S\cup T}.
\end{align*}
It remains to prove that the $2^k$ faces are distinct.  If
$S\nsubseteq T$, choose $i\in S\setminus T$.  Then
$A_i(\mathbf H)\le F_S$ and $F_T\le H_i$.  Equality $F_S=F_T$ would imply
$A_i(\mathbf H)\le H_i$, whereas
\[
 A_i(\mathbf H)\wedge H_i=X_{\mathbf H}.
\]
This contradicts Equation~\eqref{eq:atomic-rise}.  Hence
$S\mapsto F_S$ is a lattice embedding of $B_k$, and
$F_{[k]\setminus\{i\}}=H_i$.  Its image is exactly the sublattice generated
by the $H_i$ and $\hat1$, so $\mathbf H$ orders a Boolean antichain.
% relation-claim-end: proof-global-criterion

% relation-claim-begin: proof-global-enumerator
The conditions also force every $H_i$ to be proper and the entries to be
distinct: a top or repeated entry would make some $A_i$ equal to
$X_{\mathbf H}$.  Thus the symmetric group acts freely on the admissible
tuples, and every unordered Boolean antichain has exactly $k!$ orderings.
This proves Equation~\eqref{eq:global-enumerator}.
% relation-claim-end: proof-global-enumerator
% relation-claim-begin: proof-global-boundaries
The height bound follows
from the strict chain $F_\varnothing<F_{\{1\}}<\cdots<F_{[k]}$.
% relation-claim-end: proof-global-boundaries
\end{proof}

Only $2^k$ height gaps occur in Equation~\eqref{eq:global-enumerator}; the gaps
for $S=\varnothing$ and $|S|=1$ vanish identically.  Thus the criterion gives
an exact finite algorithm using $O(2^k)$ lattice operations per ordered
tuple.  This replaces the pairwise test on all $2^k$ generated faces.

The full enumerator also has an exact product law.  Write
$c_k(L)=|\cC_k(L)|$, $e_k(L)=k!c_k(L)$, and
\begin{equation}\label{eq:hom-transform}
% formal-claim-begin: def-hom-transform
 \eta_k(L)=\sum_{a=0}^k\binom ka e_a(L).
% formal-claim-end: def-hom-transform
\end{equation}

\begin{theorem}[Lattice-product convolution]\label{thm:product}
% relation-claim-begin: lead-product-transform
For finite lattices $L$ and $K$,
% relation-claim-end: lead-product-transform
\begin{equation}\label{eq:product-transform}
% formal-claim-begin: eq-product-transform
 \eta_k(L\times K)=\eta_k(L)\eta_k(K).
% formal-claim-end: eq-product-transform
\end{equation}
% relation-claim-begin: lead-product-convolution
Equivalently,
% relation-claim-end: lead-product-convolution
\begin{equation}\label{eq:product-convolution}
% formal-claim-begin: eq-product-convolution
 \boxed{
 c_k(L\times K)=
 \sum_{\substack{0\le a,b\le k\\a+b\ge k}}
 \frac{a!b!}{(k-a)!(k-b)!(a+b-k)!}\,c_a(L)c_b(K).}
% formal-claim-end: eq-product-convolution
\end{equation}
\end{theorem}

\begin{proof}
% relation-claim-begin: proof-product-transform
The quantity $e_a(L)$ counts labelled lattice embeddings $B_a\hookrightarrow
L$ that preserve the top.  A top-preserving lattice homomorphism
$\phi:B_k\to L$ may collapse some Boolean atoms to the image of the bottom.
If $A\subseteq[k]$ is the set of uncollapsed atoms, the restriction to those
atoms is a labelled embedding of $B_{|A|}$.  Indeed, if two subsets of $A$
had the same image, an atom in their symmetric difference would lie below
that common image while meeting it in the image of the bottom, forcing that
atom to be collapsed.  The collapsed atoms do not change joins, so this
embedding still preserves the top of $L$.  Conversely, a top-preserving
embedding on $A$ extends uniquely by sending a subset of $[k]$ to the image
of its intersection with $A$.  Hence $\eta_k(L)$ counts all
top-preserving homomorphisms $B_k\to L$.

A homomorphism to $L\times K$ is a pair of coordinate homomorphisms, proving
Equation~\eqref{eq:product-transform}.
% relation-claim-end: proof-product-transform
% relation-claim-begin: proof-product-convolution
The product map is injective precisely
when the two active sets $A,B\subseteq[k]$ cover $[k]$.  For
$|A|=a$, $|B|=b$, the number of such pairs is
\[
 \frac{k!}{(k-a)!(k-b)!(a+b-k)!}.
\]
Multiplying by $e_a(L)e_b(K)$, summing, and dividing by $k!$ gives
Equation~\eqref{eq:product-convolution}.
% relation-claim-end: proof-product-convolution
\end{proof}

For matroids, $\cL(M_1\oplus M_2)=\cL(M_1)\times\cL(M_2)$, so
Equation~\eqref{eq:product-convolution} computes every Boolean-antichain layer
of a direct sum from those of its summands.  Taking $K=B_1$ gives
\[
 c_k(L\times B_1)=(k+1)c_k(L)+c_{k-1}(L).
\]
Iteration from the one-element lattice recovers
$c_k(B_n)=S(n+1,k+1)$ together with the Stirling recurrence.

For $L=\cL(M)$, lattice height is matroid rank.  Thus
Theorem~\ref{thm:global-gap} specializes to a rank-gap certificate in every
finite geometric lattice.

The ambient ranks of the reconstructed faces retain additional structure.

\begin{corollary}[Rank profiles]\label{cor:rank-profile}
% formal-claim-begin: def-rank-profile
For $C\in\cC_k(M)$, put
\[
 \rho_C(S)=\rk_M(F_S)-\rk_M(X_C).
\]
% formal-claim-end: def-rank-profile
% formal-claim-begin: cor-rank-profile-polymatroid
Then $\rho_C$ is a normalized, strictly increasing integral polymatroid rank
function on $[k]$.
% formal-claim-end: cor-rank-profile-polymatroid
% formal-claim-begin: cor-rank-profile-tight
Its rank-tight members are characterized by
$\rho_C(S)=|S|$ for every $S\subseteq[k]$.
% formal-claim-end: cor-rank-profile-tight
\end{corollary}

\begin{proof}
% relation-claim-begin: proof-rank-profile-polymatroid
Normalization and strict increase follow from $F_\varnothing=X_C$ and the
strict Boolean embedding.  Matroid submodularity and
$F_S\cap F_T=F_{S\cap T}$, $F_S\vee F_T=F_{S\cup T}$ give
\[
 \rho_C(S)+\rho_C(T)\ge
 \rho_C(S\cap T)+\rho_C(S\cup T).
\]
% relation-claim-end: proof-rank-profile-polymatroid
% relation-claim-begin: proof-rank-profile-tight
If $C$ is rank-tight, every one of the $k$ strict steps in a maximal Boolean
chain raises ambient rank by one, giving $\rho_C(S)=|S|$.  The converse is
immediate at $S=[k]$.
% relation-claim-end: proof-rank-profile-tight
\end{proof}

The criterion also turns the full uniform-matroid problem into a closed
block count.

\begin{theorem}[All sizes in a uniform matroid]\label{thm:uniform-all}
% formal-claim-begin: thm-uniform-base-layers
Let $1\le r\le n$.  Then $|\cC_0(U_{r,n})|=1$ and
\[
 |\cC_1(U_{r,n})|=\sum_{b=0}^{r-1}\binom nb.
\]
% formal-claim-end: thm-uniform-base-layers
% relation-claim-begin: lead-uniform-all-formula
For $2\le k\le r$, write $p=p_1+\cdots+p_k$.  Then
% relation-claim-end: lead-uniform-all-formula
\begin{equation}\label{eq:uniform-all}
% formal-claim-begin: eq-uniform-all
 \boxed{
 |\cC_k(U_{r,n})|
 =\frac1{k!}\sum_{b=0}^{r-1}\binom nb
 \sum_{\substack{p_1,\ldots,p_k\ge1,\ p\le n-b\\
                   b+p\ge r,\ b+p-p_i<r\ (i\in[k])}}
 \frac{(n-b)!}{(n-b-p)!\prod_{i=1}^k p_i!}.}
% formal-claim-end: eq-uniform-all
\end{equation}
% formal-claim-begin: thm-uniform-distributions
In particular, the size distributions of $U_{3,4}$ and $U_{4,5}$ are
respectively
\[
 (1,11,27,4),\qquad (1,26,150,50,5).
\]
% formal-claim-end: thm-uniform-distributions
\end{theorem}

\begin{proof}
% relation-claim-begin: proof-uniform-base-layers
For $k=1$, every proper flat is a Boolean antichain, and the proper flats of
$U_{r,n}$ are the subsets of size less than $r$.
% relation-claim-end: proof-uniform-base-layers

% relation-claim-begin: proof-uniform-structure
Let $k\ge2$ and let $X$ and $A_i$ be the bottom and atoms of a Boolean span.
All are proper uniform-matroid flats.  Since $A_i\cap A_j=X$ for $i\ne j$,
the sets $P_i=A_i\setminus X$ are pairwise disjoint and nonempty.  Put
$b=|X|$ and $p_i=|P_i|$.  Every proper Boolean face is the ordinary union
\[
 X\cup\bigcup_{i\in S}P_i.
\]
The full-join and proper-coatom conditions are respectively
$b+p\ge r$ and $b+p-p_i<r$ for every $i$; these inequalities then make every
proper face a flat.  Conversely, any such bottom and blocks produce
all the faces of a Boolean sublattice, so this construction is bijective.
% relation-claim-end: proof-uniform-structure

% relation-claim-begin: proof-uniform-count
There are $\binom nb$ choices for $X$.  For fixed ordered block sizes, the
number of ordered disjoint blocks is
\[
 \frac{(n-b)!}{(n-b-p)!\prod_i p_i!}.
\]
Division by the free permutation action on the $k$ nonempty blocks gives
Equation~\eqref{eq:uniform-all}.
% relation-claim-end: proof-uniform-count
\end{proof}

\section{Two-element Boolean antichains}

The first off-diagonal layer has a closed form in every finite lattice, not
only in geometric lattices.  For a finite lattice $L$, put
\[
 \ell(X)=|[\hat0,X]|,
\]
and let $\mu_L$ be its M\"obius function.

\begin{theorem}[M\"obius formula for pairs]\label{thm:size-two}
% formal-claim-begin: thm-size-two-structural
Two elements $F,G$ form a Boolean antichain if and only if both are proper
and $F\vee G=\hat1$.
% formal-claim-end: thm-size-two-structural
% relation-claim-begin: lead-size-two-mobius
Consequently,
% relation-claim-end: lead-size-two-mobius
\begin{equation}\label{eq:size-two-mobius}
% formal-claim-begin: eq-size-two-mobius
 |\cC_2(L)|=\frac12\left(
  \sum_{X\le\hat1}\mu_L(X,\hat1)\ell(X)^2-2|L|+1
 \right).
% formal-claim-end: eq-size-two-mobius
\end{equation}
\end{theorem}

\begin{proof}
% relation-claim-begin: proof-size-two-structural
If $F,G<\hat1$ and $F\vee G=\hat1$, then neither can contain the other.  The
four elements
\[
 \hat1,\quad F,\quad G,\quad F\wedge G
\]
are distinct and form the Boolean span prescribed by the meet map.  The
converse follows because the join of the two Boolean coatoms is the top.
% relation-claim-end: proof-size-two-structural

% relation-claim-begin: proof-size-two-mobius
For $Y\in L$, let $g(Y)$ be the number of ordered pairs $(F,G)$ with
$F\vee G=Y$.  The number of ordered pairs whose join is at most $Y$ is
$\ell(Y)^2$, so
\[
 \ell(Y)^2=\sum_{X\le Y}g(X).
\]
M\"obius inversion gives
\[
 g(\hat1)=\sum_{X\le\hat1}\mu_L(X,\hat1)\ell(X)^2.
\]
Among these ordered pairs, exactly $2|L|-1$ have at least one coordinate
equal to $\hat1$.  Every remaining pair consists of incomparable proper
elements and is counted twice as an unordered Boolean antichain.  This proves
Equation~\eqref{eq:size-two-mobius}.
% relation-claim-end: proof-size-two-mobius
\end{proof}

For $L=B_n$, one has $\ell(X)=2^{\rk(X)}$ and
$\mu(X,\hat1)=(-1)^{n-\rk(X)}$.  The sum in
Equation~\eqref{eq:size-two-mobius} is therefore $3^n$, and
\[
 |\cC_2(B_n)|=\frac{3^n-2^{n+1}+1}{2}=S(n+1,3),
\]
recovering the size-two instance of the source formula.

\section{All sizes in distributive lattices}

Let $P$ be a finite poset and let $J(P)$ be its distributive lattice of order
ideals.  Denote by $\mathcal F^+(P)$ the set of nonempty order filters of
$P$, and let $\Gamma_P$ be their intersection graph: two filters are adjacent
when they have nonempty intersection.  Write
\[
 I(\Gamma_P;t)=\sum_{k\ge0}i_k(\Gamma_P)t^k
\]
for its independence polynomial.

\begin{theorem}[Distributive-lattice classification]\label{thm:distributive}
% formal-claim-begin: thm-distributive-bijection
Complementation gives a size-preserving bijection
\[
 \cC_k(J(P))
 \longleftrightarrow
 \left\{\{F_1,\ldots,F_k\}\subseteq\mathcal F^+(P):
 F_i\cap F_j=\varnothing\text{ for }i\ne j\right\}.
\]
% formal-claim-end: thm-distributive-bijection
% relation-claim-begin: lead-distributive-polynomial
Consequently,
% relation-claim-end: lead-distributive-polynomial
\begin{equation}\label{eq:distributive-independence}
% formal-claim-begin: eq-distributive-polynomial
 \sum_{k\ge0}|\cC_k(J(P))|t^k=I(\Gamma_P;t).
% formal-claim-end: eq-distributive-polynomial
\end{equation}
\end{theorem}

\begin{proof}
% relation-claim-begin: proof-distributive-criterion
A Boolean span has pairwise top joins.  Conversely, suppose that a family of
proper elements in a distributive lattice has pairwise top joins.  Finite
distributivity gives, for subsets $S,T$ of the family,
\[
 \left(\bigwedge S\right)\vee\left(\bigwedge T\right)
 =\bigwedge_{a\in S,\,b\in T}(a\vee b)
 =\bigwedge(S\cap T),
\]
because the factors with $a\ne b$ are the top and the remaining factors are
the common elements.  This is the defining meet-map identity.
% relation-claim-end: proof-distributive-criterion

% relation-claim-begin: proof-distributive-bijection
Complementation sends the proper ideals bijectively to the nonempty order
filters.  Moreover,
\[
 I_i\cup I_j=P
 \quad\Longleftrightarrow\quad
 (P\setminus I_i)\cap(P\setminus I_j)=\varnothing.
\]
Thus complementation is the displayed bijection.
% relation-claim-end: proof-distributive-bijection
% relation-claim-begin: proof-distributive-polynomial
Pairwise disjoint filters
are independent sets of $\Gamma_P$, proving
Equation~\eqref{eq:distributive-independence}.
% relation-claim-end: proof-distributive-polynomial
\end{proof}

If $P$ is an $n$-element antichain, then its nonempty filters are the
nonempty subsets of $P$.  An independent set in $\Gamma_P$ is an unordered
family of disjoint nonempty subsets; adjoining the unused elements as one
distinguished block gives
\[
 |\cC_k(B_n)|=S(n+1,k+1).
\]
If $P$ is an $n$-element chain, all nonempty filters intersect, so the only
nonzero counts are $|\cC_0|=1$ and $|\cC_1|=n$.

\section{Maximum Boolean antichains}

Let $\cC_r(M)$ denote the set of Boolean antichains of size $r$ in
$\cL(M)$, and let $\cB(\si(M))$ denote the set of bases of the
simplification.

\begin{theorem}[Maximum-antichain bijection]\label{thm:main}
% formal-claim-begin: thm-main-bijection
For every finite rank-$r$ matroid $M$, the maps
\begin{align*}
 \Phi:\cC_r(M)&\longrightarrow\cB(\si(M)),
 &C&\longmapsto\{A_1(C),\ldots,A_r(C)\},\\
 \Psi:\cB(\si(M))&\longrightarrow\cC_r(M),
 &\{A_1,\ldots,A_r\}&\longmapsto
   \left\{\bigvee_{j\ne i}A_j:i\in[r]\right\}
\end{align*}
are well defined and inverse to one another.
% formal-claim-end: thm-main-bijection
% relation-claim-begin: lead-main-count
Hence
% relation-claim-end: lead-main-count
\begin{equation}\label{eq:max-count}
% formal-claim-begin: eq-main-count
 |\cC_r(M)|=|\cB(\si(M))|=T_{\si(M)}(1,1).
% formal-claim-end: eq-main-count
\end{equation}
\end{theorem}

\begin{proof}
% relation-claim-begin: proof-main-forward
Let $C=\{H_1,\ldots,H_r\}$ lie in $\cC_r(M)$.  By
Lemma~\ref{lem:full-height}, the flats $A_i(C)$ in
Equation~\eqref{eq:boolean-atoms} have rank one.  Their join is the top of
the Boolean span, namely $E$.  Thus they are $r$ rank-one flats of total rank
$r$, so they form a basis of $\si(M)$.  This proves that $\Phi$ is well
defined.
% relation-claim-end: proof-main-forward

% relation-claim-begin: proof-main-reverse
Conversely, let $\mathcal A=\{A_1,\ldots,A_r\}$ be a basis of $\si(M)$.
Lemma~\ref{lem:basis-span} shows that its joins form a Boolean sublattice.
The displayed family defining $\Psi(\mathcal A)$ is precisely the set of
coatoms of this sublattice.  The power-set meet map is therefore an
anti-isomorphism, so the family is a Boolean antichain.
% relation-claim-end: proof-main-reverse

% relation-claim-begin: proof-main-inverses
The atoms of the Boolean span of $\Psi(\mathcal A)$ are exactly the $A_i$,
and the coatoms of the Boolean span of $C$ are exactly the $H_i$.  Therefore
$\Phi\Psi$ and $\Psi\Phi$ are the identity maps.
% relation-claim-end: proof-main-inverses
% relation-claim-begin: proof-main-count
The last equality in
Equation~\eqref{eq:max-count} is the standard basis evaluation of the Tutte
polynomial.
% relation-claim-end: proof-main-count
\end{proof}

The theorem resolves the graphic expectation of Garber et al. at matroid
scope.  Its exact form records the necessary conventions: a nonsimple matroid
is counted after simplification, and the relevant size is the matroid rank.

\section{The multivariate basis polynomial}

The bijection loses no information about the bases of the original matroid
once its rank-one flats are retained.  Define the multivariate basis
polynomial
\[
 \mathbf B_M(\mathbf x)=\sum_{B\in\cB(M)}\mathbf x^B,
 \qquad \mathbf x^B=\prod_{e\in B}x_e.
\]

\begin{theorem}[Parallel-class decomposition]\label{thm:weighted}
% relation-claim-begin: lead-weighted-polynomial
For every finite matroid $M$,
% relation-claim-end: lead-weighted-polynomial
\begin{equation}\label{eq:weighted}
% formal-claim-begin: eq-weighted-polynomial
 \boxed{
 \mathbf B_M(\mathbf x)
 =\sum_{C\in\cC_r(M)}
   \prod_{A\in\operatorname{At}(\langle C\rangle)}
   \left(\sum_{e\in A\setminus L_0}x_e\right).}
% formal-claim-end: eq-weighted-polynomial
\end{equation}
% relation-claim-begin: lead-weighted-count
In particular, if $b(M)$ is the number of bases of $M$, then
% relation-claim-end: lead-weighted-count
\begin{equation}\label{eq:weighted-count}
% formal-claim-begin: eq-weighted-count
 b(M)=\sum_{C\in\cC_r(M)}
 \prod_{A\in\operatorname{At}(\langle C\rangle)}|A\setminus L_0|.
% formal-claim-end: eq-weighted-count
\end{equation}
\end{theorem}

\begin{proof}
% relation-claim-begin: proof-weighted-polynomial
Every basis $B$ of $M$ meets exactly $r$ distinct nonloop parallel classes.
Their rank-one flats form a basis of $\si(M)$ and hence, by
Theorem~\ref{thm:main}, determine a unique $C\in\cC_r(M)$.  Conversely, for
the atoms of the Boolean span of a fixed $C$, every choice of one element
from each $A\setminus L_0$ is a basis of $M$.  Thus the bases are partitioned
by $C$, and expanding the corresponding product of linear forms gives
Equation~\eqref{eq:weighted}.
% relation-claim-end: proof-weighted-polynomial
% relation-claim-begin: proof-weighted-count
Setting every variable equal to one gives
Equation~\eqref{eq:weighted-count}.
% relation-claim-end: proof-weighted-count
\end{proof}

For a simple matroid every factor in Equation~\eqref{eq:weighted-count} is
one, reducing it to the bijective count.  The polynomial form is the correct
statement for multigraphs and other parallel extensions.

\section{Intervals, minors, and rank-tight antichains}

Let $X\le Y$ be flats of $M$.  The interval $[X,Y]$ is canonically the
lattice of flats of
\[
 N=(M|Y)/X.
\]
Indeed, the correspondence sends a flat $F$ with $X\subseteq F\subseteq Y$
to $F\setminus X$; the contraction rank identity
\[
 \rk_N(S)=\rk_M(S\cup X)-\rk_M(X)
\]
shows directly that it preserves closed sets and inclusion.  A Boolean
antichain in this interval is defined relative
to its top $Y$.  Write $d=\rk_M(Y)-\rk_M(X)$.

\begin{theorem}[Interval form]\label{thm:interval}
% formal-claim-begin: thm-interval-bijection
The Boolean antichains of size $d$ in $[X,Y]$ are naturally in bijection with
the bases of $\si((M|Y)/X)$.
% formal-claim-end: thm-interval-bijection
% relation-claim-begin: lead-interval-weighted
Moreover,
% relation-claim-end: lead-interval-weighted
\begin{equation}\label{eq:interval-weighted}
% formal-claim-begin: eq-interval-weighted
 \mathbf B_{(M|Y)/X}(\mathbf x)
 =\sum_C\prod_{A\in\operatorname{At}(\langle C\rangle)}
       \left(\sum_{e\in A\setminus X}x_e\right),
% formal-claim-end: eq-interval-weighted
\end{equation}
% relation-claim-begin: tail-interval-weighted
where $C$ runs over the size-$d$ Boolean antichains in $[X,Y]$.
% relation-claim-end: tail-interval-weighted
\end{theorem}

\begin{proof}
% relation-claim-begin: proof-interval-bijection
Apply Theorems~\ref{thm:main} and~\ref{thm:weighted} to the minor $N$ under
the flat-lattice isomorphism $\cL(N)\cong[X,Y]$.
% relation-claim-end: proof-interval-bijection
% relation-claim-begin: proof-interval-weighted
The rank-one flat of $N$
corresponding to an interval atom $A$ has parallel class $A\setminus X$,
which gives the displayed factors.
% relation-claim-end: proof-interval-weighted
\end{proof}

\begin{definition}
% formal-claim-begin: def-rank-tight
A Boolean antichain $C$ of size $k$ in $\cL(M)$ is \emph{rank-tight} if
\[
 \rk_M(X_C)=r-k.
\]
% formal-claim-end: def-rank-tight
% formal-claim-begin: def-rank-tight-equivalent
Equivalently, its Boolean span has the full height of the interval
$[X_C,E]$.
% formal-claim-end: def-rank-tight-equivalent
\end{definition}

\begin{corollary}[Rank-tight enumeration]\label{cor:rank-tight}
% formal-claim-begin: cor-rank-tight-partition
The rank-tight Boolean antichains of size $k$ are partitioned by their bottom
flat,
% formal-claim-end: cor-rank-tight-partition
% relation-claim-begin: lead-rank-tight-count
and
% relation-claim-end: lead-rank-tight-count
\begin{equation}\label{eq:rank-tight}
% formal-claim-begin: eq-rank-tight-count
 \#\{C:|C|=k,\ C\text{ rank-tight}\}
 =\sum_{\substack{X\in\cL(M)\\\rk_M(X)=r-k}}
   b\bigl(\si(M/X)\bigr).
% formal-claim-end: eq-rank-tight-count
\end{equation}
% formal-claim-begin: cor-rank-tight-fiber
For each fixed $X$, the bijection is the interval bijection of
Theorem~\ref{thm:interval} on $[X,E]$.
% formal-claim-end: cor-rank-tight-fiber
\end{corollary}

\begin{proof}
% relation-claim-begin: proof-rank-tight-fiber
If $C$ is rank-tight with bottom $X$, then it is a maximum Boolean antichain
in the rank-$k$ interval $[X,E]$.  Conversely, every maximum Boolean
antichain in that interval is rank-tight in $\cL(M)$.
% relation-claim-end: proof-rank-tight-fiber
% relation-claim-begin: proof-rank-tight-count
Apply
Theorem~\ref{thm:interval} and sum over the possible bottom flats.
% relation-claim-end: proof-rank-tight-count
\end{proof}

\section{Closed rank-tight enumerations}

Corollary~\ref{cor:rank-tight} becomes completely explicit for the Boolean
and partition lattices treated in the source paper.

\begin{corollary}[Boolean lattices]\label{cor:boolean-lattice}
% relation-claim-begin: lead-boolean-tight
The number of rank-tight Boolean antichains of size $k$ in the Boolean
lattice $B_n$ is
% relation-claim-end: lead-boolean-tight
\begin{equation}\label{eq:boolean-tight}
% formal-claim-begin: eq-boolean-tight
 \binom nk.
% formal-claim-end: eq-boolean-tight
\end{equation}
% formal-claim-begin: cor-boolean-total
Since the total number of size-$k$ Boolean antichains is
$S(n+1,k+1)$ \cite{GarberEtAl2026},
% formal-claim-end: cor-boolean-total
% relation-claim-begin: lead-boolean-nontight
the number of non-rank-tight ones is
% relation-claim-end: lead-boolean-nontight
\begin{equation}\label{eq:boolean-nontight}
% formal-claim-begin: eq-boolean-nontight
 S(n+1,k+1)-\binom nk.
% formal-claim-end: eq-boolean-nontight
\end{equation}
\end{corollary}

\begin{proof}
% relation-claim-begin: proof-boolean-tight
The lattice $B_n$ is the flat lattice of the free matroid on $n$ elements.
A corank-$k$ flat is the complement of a $k$-subset, and the contraction by
that flat is the free matroid on those $k$ elements.  Its simplification has
one basis.  Equation~\eqref{eq:boolean-tight} follows from
Corollary~\ref{cor:rank-tight};
% relation-claim-end: proof-boolean-tight
% relation-claim-begin: proof-boolean-nontight
subtract it from the total count for the last
assertion.
% relation-claim-end: proof-boolean-nontight
\end{proof}

\begin{corollary}[Partition lattices]\label{cor:partition-lattice}
% relation-claim-begin: lead-partition-tight
For $1\le k\le n-1$, the number of rank-tight Boolean antichains of size $k$
in the partition lattice $\Pi_n$ is
% relation-claim-end: lead-partition-tight
\begin{equation}\label{eq:partition-tight}
% formal-claim-begin: eq-partition-tight
 S(n,k+1)(k+1)^{k-1}.
% formal-claim-end: eq-partition-tight
\end{equation}
% formal-claim-begin: cor-partition-maximum
At $k=n-1$, this is Cayley's value $n^{n-2}$ for maximum Boolean
antichains.
% formal-claim-end: cor-partition-maximum
\end{corollary}

\begin{proof}
% relation-claim-begin: proof-partition-structure
Identify $\Pi_n$ with $\cL(M(K_n))$.  A corank-$k$ flat is a set partition
of $[n]$ into $k+1$ blocks, of which there are $S(n,k+1)$.  Contracting the
flat collapses each block to a vertex.  Between every two resulting vertices
there is a nonempty parallel class, so the simplification of the contraction
is $M(K_{k+1})$.
% relation-claim-end: proof-partition-structure
% relation-claim-begin: proof-partition-count
It has $(k+1)^{k-1}$ bases by Cayley's formula.  Apply
Corollary~\ref{cor:rank-tight}.
% relation-claim-end: proof-partition-count
\end{proof}

The same calculation gives all rank-tight layers of a uniform matroid.

\begin{corollary}[Uniform matroids]\label{cor:uniform-tight}
% formal-claim-begin: cor-uniform-tight
For $2\le r\le n$, the number $R_{r,n}(k)$ of rank-tight size-$k$ Boolean
antichains in $\cL(U_{r,n})$ is
\[
 R_{r,n}(k)=
 \begin{cases}
 1,&k=0,\\
 \binom{n}{r-1},&k=1,\\
 \binom nr\binom rk,&2\le k\le r.
 \end{cases}
\]
% formal-claim-end: cor-uniform-tight
\end{corollary}

\begin{proof}
% relation-claim-begin: proof-uniform-tight
A corank-$k$ flat has size $r-k$.  For $k=1$, the rank-one contraction
simplifies to a one-element matroid.  For $k\ge2$, the contraction is
$U_{k,n-r+k}$, which is simple and has $\binom{n-r+k}{k}$ bases.  Thus
Corollary~\ref{cor:rank-tight} gives
\[
 \binom{n}{r-k}\binom{n-r+k}{k}
 =\binom nr\binom rk.
\]
% relation-claim-end: proof-uniform-tight
\end{proof}

The same M\"obius formula gives a compact expression for the entire
size-two layer of the partition lattice, including the non-rank-tight part.

\begin{corollary}[All Boolean pairs in $\Pi_n$]\label{cor:partition-pairs}
% formal-claim-begin: def-bell-square-egf
Let $\mathsf B_m$ be the $m$th Bell number and put
\[
 A(z)=\sum_{m\ge1}\mathsf B_m^2\frac{z^m}{m!}.
\]
% formal-claim-end: def-bell-square-egf
% relation-claim-begin: lead-partition-pairs
Then, for $n\ge1$,
% relation-claim-end: lead-partition-pairs
\begin{equation}\label{eq:partition-pairs}
% formal-claim-begin: eq-partition-pairs
 |\cC_2(\Pi_n)|=\frac12\left(
 n![z^n]\log(1+A(z))-2\mathsf B_n+1
 \right).
% formal-claim-end: eq-partition-pairs
\end{equation}
% formal-claim-begin: cor-partition-pairs-initial
The initial values for $n=1,2,\ldots$ are
\[
 0,\ 0,\ 3,\ 45,\ 620,\ 9750,\ 183680,\ldots.
\]
% formal-claim-end: cor-partition-pairs-initial
\end{corollary}

\begin{proof}
% relation-claim-begin: proof-partition-pairs-lower
For a set partition $\pi$, the lower interval is the product of the partition
lattices on its blocks, so
\[
 \ell(\pi)=\prod_{D\in\pi}\mathsf B_{|D|}.
\]
% relation-claim-end: proof-partition-pairs-lower
% relation-claim-begin: proof-partition-pairs-log
If $b(\pi)$ is the number of blocks, then
$\mu(\pi,\hat1)=(-1)^{b(\pi)-1}(b(\pi)-1)!$.  The exponential formula now
gives
\[
 \sum_{\pi\in\Pi_n}\mu(\pi,\hat1)\ell(\pi)^2
 =n![z^n]\sum_{b\ge1}\frac{(-1)^{b-1}}b A(z)^b
 =n![z^n]\log(1+A(z)).
\]
% relation-claim-end: proof-partition-pairs-log
% relation-claim-begin: proof-partition-pairs-substitute
Substitute this and $|\Pi_n|=\mathsf B_n$ into
Equation~\eqref{eq:size-two-mobius}.
% relation-claim-end: proof-partition-pairs-substitute
\end{proof}

\section{All sizes in subspace lattices}

Let $V=\mathbb F_q^n$, where $q$ is a prime power, and let $\mathcal S(V)$ be
its lattice of linear subspaces.  Write $\qbinom{n}{d}$ for the Gaussian
binomial coefficient and
\[
 g_d(q)=|\operatorname{GL}_d(q)|
       =\prod_{i=0}^{d-1}(q^d-q^i),
 \qquad g_0(q)=1.
\]
For $1\le k\le d$, put
\begin{equation}\label{eq:direct-sum-count}
% formal-claim-begin: def-direct-sum-count
 D_{d,k}(q)=\frac1{k!}
 \sum_{\substack{d_1+\cdots+d_k=d\\d_i\ge1}}
 \frac{g_d(q)}{\prod_{i=1}^k g_{d_i}(q)}.
% formal-claim-end: def-direct-sum-count
\end{equation}

\begin{theorem}[Subspace-lattice enumeration]\label{thm:subspace}
% formal-claim-begin: thm-subspace-correspondence
Every size-$k$ Boolean antichain in $\mathcal S(V)$ corresponds uniquely to:

\begin{enumerate}
\item a subspace $X\le V$ of codimension $d\ge k$;
\item an unordered direct-sum decomposition
\[
 V/X=U_1\oplus\cdots\oplus U_k
\]
into nonzero subspaces.
\end{enumerate}
% formal-claim-end: thm-subspace-correspondence

% relation-claim-begin: lead-subspace-count
Hence, for $k\ge1$,
% relation-claim-end: lead-subspace-count
\begin{equation}\label{eq:subspace-count}
% formal-claim-begin: eq-subspace-count
 \boxed{
 |\cC_k(\mathcal S(\mathbb F_q^n))|
 =\sum_{d=k}^n \qbinom{n}{d} D_{d,k}(q).}
% formal-claim-end: eq-subspace-count
\end{equation}
% formal-claim-begin: thm-subspace-zero-layer
Together with $|\cC_0|=1$, this enumerates all Boolean antichains in every
finite subspace lattice.
% formal-claim-end: thm-subspace-zero-layer
\end{theorem}

\begin{proof}
% relation-claim-begin: proof-subspace-correspondence
Let $C=\{H_1,\ldots,H_k\}$ be Boolean, put $X=\bigcap_iH_i$, and let
$A_i=\bigcap_{j\ne i}H_j$ be the atoms of its Boolean span.  In the quotient
$V/X$, set $U_i=A_i/X$.  The Boolean identities give
\[
 \sum_iU_i=V/X,
 \qquad
 U_i\cap\sum_{j\ne i}U_j=0,
\]
so the sum is direct and every $U_i$ is nonzero.  Conversely, from such a
decomposition define
\[
 H_i/X=\bigoplus_{j\ne i}U_j.
\]
The sums of subcollections of the $U_i$ form a Boolean sublattice, whose
coatoms are the $H_i$.  The two constructions recover $X$, the summands, and
the coatoms, so they are inverse.
% relation-claim-end: proof-subspace-correspondence

% relation-claim-begin: proof-subspace-count
There are $\qbinom{n}{d}$ choices for $X$.  For fixed positive dimensions
$d_1,\ldots,d_k$, the group $\operatorname{GL}_d(q)$ acts transitively on
ordered direct-sum decompositions with stabilizer
$\prod_i\operatorname{GL}_{d_i}(q)$.  Summing over ordered positive
compositions and dividing by $k!$ gives Equation~\eqref{eq:direct-sum-count},
then Equation~\eqref{eq:subspace-count}.
% relation-claim-end: proof-subspace-count
\end{proof}

For example, the size distribution for the subspace lattice of
$\mathbb F_2^3$ is
\[
 1,\ 15,\ 49,\ 28.
\]
At $k=n$, Equation~\eqref{eq:subspace-count} reduces to the projective-basis
formula in Equation~\eqref{eq:projective-count} below.

\section{Graphic and classical specializations}

Let $G$ be a finite graph, allowing loops and parallel edges, and let $c(G)$
be its number of connected components.  Its cycle matroid has rank
$|V(G)|-c(G)$.  Simplifying the cycle matroid deletes loops and identifies
parallel edges.

\begin{corollary}[Spanning forests]\label{cor:graphic}
% formal-claim-begin: cor-graphic-simplified-forests
The maximum Boolean antichains in $\cL(M(G))$ are in bijection with the
spanning forests of the simplification of $G$.
% formal-claim-end: cor-graphic-simplified-forests
% formal-claim-begin: cor-graphic-size-specializations
Their size is
$|V(G)|-c(G)$.  If $G$ is simple, they are in bijection with the spanning
forests of $G$; if $G$ is also connected, they are in bijection with its
spanning trees.
% formal-claim-end: cor-graphic-size-specializations

% formal-claim-begin: cor-graphic-weighted
With an edge variable $x_e$, Equation~\eqref{eq:weighted} is the spanning-
forest polynomial of a multigraph: each simplified edge contributes the sum
of the variables in its parallel class.
% formal-claim-end: cor-graphic-weighted
\end{corollary}

\begin{proof}
% relation-claim-begin: proof-graphic
The bases of a cycle matroid are the spanning forests having one spanning
tree in each connected component.  Apply Theorems~\ref{thm:main}
and~\ref{thm:weighted}.
% relation-claim-end: proof-graphic
\end{proof}

Taking $G=K_n$, the flat lattice is the partition lattice and Cayley's formula
gives $n^{n-2}$ maximum Boolean antichains, recovering the result of
Garber et al.\ \cite{GarberEtAl2026}.

Every rank-three flat lattice is already completely enumerated by the
one-element, pair, and maximum theorems:
\begin{align*}
 c_0(M)&=1,& c_1(M)&=|\cL(M)|-1,\\
 c_2(M)&=\frac12\left(\sum_{X\le E}\mu(X,E)\ell(X)^2
              -2|\cL(M)|+1\right),&
 c_3(M)&=b(\si(M)).
\end{align*}
All higher coefficients vanish.  In particular, let $\mathsf P(q)$ be the
flat lattice of any finite projective plane of order $q$, and put
$N=q^2+q+1$.  Counting pairs of lines, nonincident point-line pairs, and
noncollinear point triples gives the complete distribution
\begin{equation}\label{eq:projective-plane-all}
% formal-claim-begin: eq-projective-plane-all
 \boxed{
 (c_0,c_1,c_2,c_3)=
 \left(1,\ 2N+1,\ \binom N2+Nq^2,\
 \binom N3-N\binom{q+1}3\right).}
% formal-claim-end: eq-projective-plane-all
\end{equation}
This calculation uses only the incidence axioms and therefore includes
non-Desarguesian finite projective planes.

For $r\ge1$ and the projective geometry $PG(r-1,q)$, an unordered basis is an unordered
set of $r$ independent projective points.  Hence
\begin{equation}\label{eq:projective-count}
% formal-claim-begin: eq-projective-count
 |\cC_r(PG(r-1,q))|
 =\frac1{r!}\prod_{i=0}^{r-1}\frac{q^r-q^i}{q-1},
 \qquad r\ge1.
% formal-claim-end: eq-projective-count
\end{equation}

\section{Scope of the matroid correspondences}

The global classification is unconditional for finite lattices.  Its matroid
correspondences take particularly precise forms.

\paragraph{Parallel classes.}
The unweighted maximum-antichain correspondence naturally counts bases of the
simplification.  For a rank-two matroid with parallel classes
$\{a,a'\}$ and $\{b,b'\}$, the flat lattice has one maximum Boolean
antichain, whereas the original matroid has four bases.  The weighted formula
recovers the latter count as $4=2\cdot2$.

\paragraph{Graphic specialization.}
For a disconnected graph, the relevant size is the matroid rank
$|V(G)|-c(G)$ and the corresponding objects are spanning forests.
Connected simple graphs specialize this statement to spanning trees, as in
Corollary~\ref{cor:graphic}.

\medskip\noindent\textbf{Off-diagonal layers.}\par\noindent
Contraction bases characterize the rank-tight diagonal.  The global
classification also includes smaller non-rank-tight Boolean antichains.  In
$U_{3,4}$, for example, the two flats
\[
 \{1,2\},\qquad\{3,4\}
\]
form a Boolean antichain of size two: their meet is empty and their join is
the top.  Its bottom has rank zero, whereas rank-tightness would require rank
one.  Theorem~\ref{thm:global-gap} detects this span through its
atom-reconstruction gaps and, more generally, enumerates every off-diagonal
layer.

\section{Further directions}

Equation~\eqref{eq:global-enumerator} determines the complete two-parameter
enumerator by antichain size and bottom rank.  The rank-tight diagonal reduces
to contraction bases, while the off-diagonal terms record realized strict
polymatroid profiles.  A natural next question is whether substantial parts
of this profile distribution factor through familiar matroid invariants, or
whether it supplies explicit separations beyond the Tutte polynomial.

The source paper also asks about the partial order on Boolean antichains.
Under Theorem~\ref{thm:main}, any natural order on maximum Boolean antichains
becomes an order on simplified matroid bases.  Determining whether it agrees
with a known basis order for special matroid classes opens a further
structural direction.

\begingroup
\footnotesize
\medskip
\noindent
% scope-claim-begin: scope-formalization-status
\textbf{Formalization.}
The accompanying Lean~4 development derives every formally mapped result from
mathlib under a kernel-only trust boundary and has passed per-declaration
axiom audits.
% scope-claim-end: scope-formalization-status

\medskip
\noindent\textbf{Code and formalization availability.}\par
The Lean~4 sources, theorem map, verification instructions, and finite
regression scripts are available at
\href{https://github.com/crabsatellite/matroid-boolean-antichains}
{\nolinkurl{github.com/crabsatellite/matroid-boolean-antichains}}.

\medskip
\noindent\textbf{Declaration of Generative AI and AI-assisted technologies in
the writing process.}\par
During the preparation of this work, the author used an author-developed local
implementation of the Proof Engine framework for AI-assisted mathematical
research.  The implementation operates through replaceable AI-agent and
harness components; the present work used OpenAI Codex within that
implementation.  The framework is described in \emph{Proof Engine
Infrastructure}~\cite{LiProofEngine2026}.  The author reviewed and edited the
resulting material and assumes full responsibility for the definitions,
statements, proofs, citations, code, and exposition.
\endgroup

\begingroup
\small
\bibliographystyle{plain}
\bibliography{references}
\endgroup
\end{document}
